\documentclass[10pt,a4paper]{amsart}
\usepackage[margin=19mm]{geometry}
\usepackage{amsmath,amssymb,mathtools}
\usepackage[scr=rsfs,cal=euler]{mathalfa}
\usepackage{xfrac}
\usepackage[unicode,hidelinks,bookmarks=false]{hyperref}
\newcommand{\ie}{i.e.\ }
\newcommand{\cf}{cf.\ }
\newcommand{\resp}{respectively\ }
\newcommand{\Subsection}{\subsection}
\providecommand{\nobreakdash}{\nobreak\hbox{-}}
\setlength{\emergencystretch}{2em}
\multlinegap0pt
\usepackage{bm}
\usepackage{bbm}
\usepackage{dsfont}
\usepackage{xspace}
\usepackage{cancel}
\usepackage{mathtools}
\usepackage{amsthm}
\makeatletter
\@ifundefined{lem}{\newtheorem{lem}{Lemma}}{}
\makeatother
\def\cR{\mathcal{R}}
\def\signn{{n(n+1)/2}}
\title[The cyclic open--closed map and variations of Hodge structur]{The cyclic open--closed map and variations of Hodge structures}
\author{Sheel Ganatra, Nick Sheridan}
\date{Source: arXiv:2511.04498v1 (CC BY 4.0)}
\begin{document}
\maketitle

\noindent\emph{Source and licence.} The cyclic open--closed map and variations of Hodge structures. Version arXiv:2511.04498v1 (CC BY 4.0), distributed under the Creative Commons Attribution 4.0 International licence, as recorded in the arXiv licence metadata for this exact version and shown on the abstract page https://arxiv.org/abs/2511.04498v1. Full rights notice: licenses/arxiv-2511.04498-CC-BY-4.0.txt.\par\smallskip

\noindent\textit{Editorial scope.} Counted unit: the Lem whose source label is \texttt{None} in the article (source0.tex lines 1638--1646, source label \texttt{None}), delivered together with the article's own complete original proof of it (source0.tex lines 1647--1661, one \texttt{proof} environment of 15 lines). No mathematics is written, repaired or re-derived: statement and proof are the article's own text line for line. Required context retained uncounted: eq:H12 vee vee (lines 1297--1299); eq:RH3 or (lines 1616--1618). Every definition, display and label of those lines is retained unchanged. Omitted from this package and not redistributed: 1--1296, 1300--1615, 1619--1637, 1662--1698. Nothing inside a retained excerpt is omitted or altered: every retained body is delivered exactly as the article's lines are written, so no omission receipt of any kind accompanies this package. Citations: the retained excerpts carry no citation command at all, so no bibliography entry of the article is transcribed and no reference is redistributed. Reference adaptation: none was needed and none was made; every reference inside the delivered unit resolves inside it, so no unresolved reference or citation is produced. Printed slips kept literal: no printed word, symbol, hypothesis or number of the counted statement or of its proof was altered. Layout and class: the article's own class is \texttt{article} with packages including amsmath, amssymb, amsthm, tikz-cd, enumerate, mathtools, appendix; the wrapper is the lane's pinned \texttt{amsart} editorial wrapper at 10pt on A4, which restates the theorem-like environments the retained text needs and adds the article's own macro definitions that the retained text needs (\texttt{\textbackslash cR}, \texttt{\textbackslash signn}), with extra wrapper packages (bm, bbm, dsfont, xspace, cancel, mathtools). The delivered numbering is the unit's own. Assets: no retained span contains a figure environment, a graphics-inclusion command, a PDF-page inclusion, a table or a TikZ picture; no image file is copied into this package, the package declares no asset dependency and no image file is redistributed with it. Licence: version arXiv:2511.04498v1 of this article is distributed under the Creative Commons Attribution 4.0 International licence, as recorded in the arXiv licence metadata for this exact version and shown on the abstract page https://arxiv.org/abs/2511.04498v1. A full rights notice is delivered at licenses/arxiv-2511.04498-CC-BY-4.0.txt. Characters: every retained excerpt is pure ASCII, so the delivered engine, XeLaTeX with its default Latin Modern text font, reports no missing character.
    \begin{multline}\label{eq:H12 vee vee}
        H^{12}_{\langle k^\vee_{in},k^\vee_{out}\rangle}(\alpha,\beta) - (-1)^\signn\delta_{k_{in},0}\delta_{k_{out},0} \langle \alpha,\beta \rangle_{Muk,\vee\vee} = \partial(H^2_{\langle k^\vee_{in},k^\vee_{out}\rangle}(\alpha,\beta)) \\+ H^2_{\langle k^\vee_{in},k^\vee_{out}\rangle}(b\alpha,\beta) + H^2_{\langle k^\vee_{in},k^\vee_{out}\rangle}(\alpha,b\beta) + H^2_{\langle (k_{in}-1)^\wedge,k^\vee_{out}\rangle}(B\alpha,\beta) \\+ H^2_{\langle k^\vee_{in},(k_{out}-1)^\wedge\rangle}(\alpha,B\beta) - H^3_{\langle (k_{in}-1)^\vee,k^\vee_{out}\rangle}(\alpha,\beta) - H^3_{\langle k^\vee_{in},(k_{out}-1)^\vee\rangle}(\alpha,\beta)
    \end{multline}

\begin{equation}\label{eq:RH3 or}
    \sigma(\cR(H^3_{\langle k_{in}^\vee,k_{out}^\vee\rangle},(1,1),(0,0))) = \sigma(R)\sigma(\theta)\sigma(2k_{in}+2k_{out}),
\end{equation}

\begin{lem}
    The signs of the isomorphisms
    \begin{align*}
    \S(H^2_{k_{in}^\vee,k_{out}^\vee}) & \cong \sigma(\partial)^\vee \S(H^3_{(k_{in}-1)^\vee,k_{out}^\vee}),\\
    \S(H^2_{k_{in}^\vee,k_{out}^\vee}) &\cong \sigma(\partial)^\vee \S(H^3_{k_{in}^\vee,(k_{out}-1)^\vee}),
    \end{align*}
    induced at the codimension-$1$ boundary component $\{z^{in}_{k_{in}}=0\}$ (respectively $\{z^{out}_{k_{out}} = 0\}$), are both $-1$. 
    This justifies the signs in front of the corresponding terms in \eqref{eq:H12 vee vee}.
\end{lem}

\begin{proof}
    In a neighbourhood of such a boundary component, let $z^{in}_{k_{in}} = r^{in} \exp( i\theta^{in})$. 
    Then the identification
    $$\sigma(\cR(H^3)) \sigma(\partial) \cong \sigma(\cR(H^2))$$
    is induced by 
    \begin{align*}
        \sigma(2) \sigma(\partial) &\cong \sigma(R)\sigma(\theta)\sigma(\partial) \\
        & \cong \sigma(R)\sigma(\theta)\sigma(r_{in}) \\
        & \cong \sigma(R)\sigma(r_{in})\sigma(\theta_{in}) \\
        &\cong \sigma(\partial) \sigma(2),
    \end{align*}
    where the first isomorphism is the isomorphism introduced in the definition of $\S(H^3)$; the second arises from the identification $\sigma(\partial) = \sigma(r_{in})$ at this boundary component; the third arises from the identification $\sigma(\theta) = \sigma(N_{H^3/H^2}) = \sigma(\theta_{in})$ introduced in the definition of $\S(H^2)$; and the final one arises from the complex orientation on the $z^{in}_{k_{in}}$-plane, which goes into the isomorphism \eqref{eq:RH3 or}, together with the identification $\sigma(R) = \sigma(\partial)$, which goes into the definition of $\S(H^2)$.
    The overall sign is easily verified to be $-1$, arising from commuting $\sigma(\theta_{in})$ with $\sigma(r_{in})$. 
    The computation of the second sign is identical.
\end{proof}

\end{document}
